% !TeX root = ../../Book4.tex
% 纲要标题：### 1.2 函子与自然变换
% 纲要位置：计划纲要.md，第 2457 行。

\section{函子与自然变换}
\label{b4:sec:0102}

一种构造若要适用于全部同态，就必须同时说明它怎样改变对象和映射。比较两种这样的构造时，所选分量还须与每个原同态交换。

% 纲要内容：
%
% * 协变、反变及自然性方块
% * 函子复合与自然变换复合
% * 忘却结构和附加结构的例子
% * 生成态射上的自然性检验；逐对象同构但不自然同构的两对象反例
%

\subsection{协变函子、反变函子与自然性方块}
\label{b4:subsec:0102:01}

把一个模忘掉标量作用后变成交换群，是一种对对象的操作；把模同态也看成群同态，才说明这种操作怎样与映射相容。函子同时规定这两部分。只列出每个对象要变成什么，还没有给出函子。

\begin{definition}\label{b4:def:functor}
设 \(\mathcal C,\mathcal D\) 为范畴。给每个 \(X\in\mathcal C\) 指定对象 \(F(X)\in\mathcal D\)，给每个态射 \(f:X\rightarrow Y\) 指定态射
\(F(f):F(X)\rightarrow F(Y)\)。若
\[
F(1_X)=1_{F(X)},\qquad F(g\circ f)=F(g)\circ F(f)
\]
对所有对象及可复合态射成立，称这组指定为从 \(\mathcal C\) 到 \(\mathcal D\) 的\term[xiebianhanzi]{协变函子}[covariant functor]，记为 \(F:\mathcal C\rightarrow\mathcal D\)。本书简称“函子”时指协变函子。
\end{definition}

\begin{definition}\label{b4:def:contravariant-functor}
设 \(\mathcal C,\mathcal D\) 为范畴。给每个对象 \(X\in\mathcal C\) 指定对象 \(F(X)\in\mathcal D\)，给每个 \(\mathcal C\) 中的态射 \(f:X\rightarrow Y\) 指定反向态射
\(F(f):F(Y)\rightarrow F(X)\)。若
\[
F(1_X)=1_{F(X)},\qquad F(g\circ f)=F(f)\circ F(g),
\]
对所有对象及可复合态射成立，称这组指定为从 \(\mathcal C\) 到 \(\mathcal D\) 的\term[fanbianhanzi]{反变函子}[contravariant functor]。
\end{definition}

写出一个具体函子时，要同时说明它把对象送到什么对象、把态射送到什么态射。下面以两个 Hom 函子为例，分别写出这两种对应。

\ProseParagraphBreak
固定局部小范畴 \(\mathcal C\) 中的对象 \(A\)。从 \(A\) 到每个对象 \(X\) 的全体态射组成集合 \(\operatorname{Hom}_{\mathcal C}(A,X)\)。先写出协变 Hom 函子在对象上的对应：
\[
\begin{aligned}
\operatorname{Hom}_{\mathcal C}(A,-):\mathcal C&\longrightarrow\mathbf{Set},\\
X&\longmapsto\operatorname{Hom}_{\mathcal C}(A,X).
\end{aligned}
\]

再给定态射 \(f\in\operatorname{Hom}_{\mathcal C}(X,Y)\)。它在这个函子作用下对应到下面的集合映射：
\[
\begin{aligned}
\operatorname{Hom}_{\mathcal C}(A,-)(f):
\operatorname{Hom}_{\mathcal C}(A,X)&\longrightarrow
\operatorname{Hom}_{\mathcal C}(A,Y),\\
u&\longmapsto f\circ u.
\end{aligned}
\]

因此，固定两个对象 \(X,Y\)，该函子在它们之间的态射集上的作用可以写为
\[
\begin{aligned}
\operatorname{Hom}_{\mathcal C}(A,-):
\operatorname{Hom}_{\mathcal C}(X,Y)&\longrightarrow
\operatorname{Hom}_{\mathbf{Set}}\bigl(
\operatorname{Hom}_{\mathcal C}(A,X),
\operatorname{Hom}_{\mathcal C}(A,Y)\bigr),\\
f&\longmapsto\bigl(u\longmapsto f\circ u\bigr).
\end{aligned}
\]
这里右端是两个集合之间的全体映射；左端的每个态射 \(f\) 被送到其中一个映射，而 \(u\) 是这个映射的输入。

\ProseParagraphBreak
符号中的横线 \(-\) 表示随对象变化的位置，\(A\) 始终固定。这里能以 \(\mathbf{Set}\) 为目标范畴，正是因为 \(\mathcal C\) 局部小，每个 Hom 类都是集合。

\ProseParagraphBreak
另一种做法是取从 \(X\) 到 \(A\) 的全体态射。相应的反变 Hom 函子在对象上的对应为
\[
\begin{aligned}
\operatorname{Hom}_{\mathcal C}(-,A):\mathcal C&\longrightarrow\mathbf{Set}
\quad\text{（反变）},\\
X&\longmapsto\operatorname{Hom}_{\mathcal C}(X,A).
\end{aligned}
\]

给定态射 \(f\in\operatorname{Hom}_{\mathcal C}(X,Y)\)，它对应到的集合映射为
\[
\begin{aligned}
\operatorname{Hom}_{\mathcal C}(-,A)(f):
\operatorname{Hom}_{\mathcal C}(Y,A)&\longrightarrow
\operatorname{Hom}_{\mathcal C}(X,A),\\
v&\longmapsto v\circ f.
\end{aligned}
\]
它把从 \(Y\) 到 \(A\) 的态射 \(v\) 送到从 \(X\) 到 \(A\) 的态射 \(v\circ f\)，所以映射的方向与 \(f\) 相反。

\ProseParagraphBreak
于是，对固定的 \(X,Y\)，该函子在态射集上的作用为
\[
\begin{aligned}
\operatorname{Hom}_{\mathcal C}(-,A):
\operatorname{Hom}_{\mathcal C}(X,Y)&\longrightarrow
\operatorname{Hom}_{\mathbf{Set}}\bigl(
\operatorname{Hom}_{\mathcal C}(Y,A),
\operatorname{Hom}_{\mathcal C}(X,A)\bigr),\\
f&\longmapsto\bigl(v\longmapsto v\circ f\bigr).
\end{aligned}
\]

恒等律分别来自 \(1_X\circ u=u\)、\(v\circ1_X=v\)，复合规律分别来自
\[
(g\circ f)\circ u=g\circ(f\circ u),\qquad
v\circ(g\circ f)=(v\circ g)\circ f.
\]
这类 Hom 函子的系统作用将在第3章讨论。

\ProseParagraphBreak
域 \(k\) 上的代数对偶也是反变函子：\(V^*=\operatorname{Hom}_k(V,k)\)，而
\(f^*(\lambda)=\lambda f\)。若 \(V\xrightarrow{f}W\xrightarrow{g}Z\)，则对 \(\mu\in Z^*\)，
\((gf)^*(\mu)=\mu gf=f^*(g^*(\mu))\)。第二卷的转置矩阵规则正是这个复合次序在坐标中的表达。
\symidx{\(V^*\)：向量空间的代数对偶}

\begin{proposition}\label{b4:prop:functors-preserve-isomorphisms}
设 \(\mathcal C,\mathcal D\) 为范畴，\(F:\mathcal C\rightarrow\mathcal D\) 为函子。若 \(f:X\rightarrow Y\) 是 \(\mathcal C\) 中的同构，则 \(F(f)\) 是 \(\mathcal D\) 中的同构，且
\(F(f^{-1})=F(f)^{-1}\)。反变函子也保留同构。
\end{proposition}
\begin{proof}
对等式 \(f^{-1}f=1_X\)、\(ff^{-1}=1_Y\) 应用协变函子，得到 \(F(f^{-1})\) 的左右逆等式。反变情形把乘积次序调换，仍然得到这两个逆等式。
\end{proof}

\begin{definition}\label{b4:def:natural-transformation}
设 \(\mathcal C,\mathcal D\) 为范畴，\(F,G:\mathcal C\rightarrow\mathcal D\) 为两个函子。对每个对象 \(X\in\mathcal C\)，给定 \(\mathcal D\) 中的态射
\(\eta_X:F(X)\rightarrow G(X)\)。若对每个 \(f:X\rightarrow Y\)，都有
\[
G(f)\eta_X=\eta_YF(f),
\]
则称这族态射为从 \(F\) 到 \(G\) 的\term[ziranbianhuan]{自然变换}[natural transformation]，记为 \(\eta:F\Rightarrow G\)，称 \(\eta_X\) 为在 \(X\) 处的分量。
\end{definition}
\symidx{\(\eta:F\Rightarrow G\)：函子之间的自然变换}

等式可以画为自然性方块：
\[
\begin{tikzcd}
F(X)\arrow[r,"\eta_X"]\arrow[d,"F(f)"']&
G(X)\arrow[d,"G(f)"]\\
F(Y)\arrow[r,"\eta_Y"']&G(Y).
\end{tikzcd}
\]
它比较两种次序：先用分量映射再沿 \(f\) 变化，或者先沿 \(f\) 变化再用分量映射。自然性要求对源范畴中的全部态射成立，但证明时不一定要逐一检验。若给定的一族态射生成源范畴，即每个态射都是这族态射的有限复合或恒等态射，则只需在生成态射上检验。事实上，恒等态射处的自然性自动成立；若分量族对 \(f:X\rightarrow Y\)、\(g:Y\rightarrow Z\) 都满足自然性，则
\[
\begin{aligned}
G(gf)\eta_X
&=G(g)G(f)\eta_X
=G(g)\eta_YF(f)\\
&=\eta_ZF(g)F(f)
=\eta_ZF(gf).
\end{aligned}
\]
因此自然性也对它们的复合成立。任意抽选若干态射，或只检验同构，通常不足以推出自然性。

\ProseParagraphBreak
两个反变函子之间的自然变换使用同样的要求，但两个函子所给的态射方向都与原箭头 \(f:X\rightarrow Y\) 相反，因此等式为
\(G(f)\eta_Y=\eta_XF(f)\)。普通自然变换要求两个函子有相同的源范畴和目标范畴。\secref{b4:sec:0103}将把反变函子表述为从对偶范畴 \(\mathcal C^{\mathrm{op}}\) 到 \(\mathcal D\) 的协变函子；要把它与 \(\mathcal C\rightarrow\mathcal D\) 比较，必须先明确共同的源范畴，不能仅因二者的对象取值相同就直接套用以上方块。

\[
\begin{tikzcd}[column sep=large,row sep=large]
F(Y)\arrow[r,"\eta_Y"]\arrow[d,"F(f)"']&
G(Y)\arrow[d,"G(f)"]\\
F(X)\arrow[r,"\eta_X"']&G(X).
\end{tikzcd}
\]

\begin{example}\label{b4:ex:torsion-natural-inclusion}
对交换群 \(A\)，记 \(t(A)=\{a:\;\text{存在整数}\;n\ge1,\ na=0\}\)。
若 \(na=0\)、\(mb=0\)，则 \(nm(a-b)=0\)，所以这是子群。群同态 \(f:A\rightarrow B\) 把挠元送到挠元，因为 \(nf(a)=f(na)=0\)。于是限制映射定义函子 \(t:\mathbf{Ab}\rightarrow\mathbf{Ab}\)，包含
\(j_A:t(A)\hookrightarrow A\) 组成自然变换 \(j:t\Rightarrow\operatorname{Id}\)。
\end{example}

这项自然性没有说挠子群都相同，也没有选取任何补群。它说每个同态都保留这个已经由对象自身确定的子群。同样，对固定整数 \(n\)，逐对象取乘 \(n\) 的同态，得到
\(\operatorname{Id}_{\mathbf{Ab}}\Rightarrow\operatorname{Id}_{\mathbf{Ab}}\)；这里自然性就是 \(f(na)=nf(a)\)。

\subsection{函子复合与自然变换复合}
\label{b4:subsec:0102:02}

函子 \(F:\mathcal C\rightarrow\mathcal D\)、\(H:\mathcal D\rightarrow\mathcal E\) 可以复合，规定
\((HF)(X)=H(F(X))\)、\((HF)(f)=H(F(f))\)。两次使用函子公理，便有
\[
HF(gf)=H\bigl(F(g)F(f)\bigr)=HF(g)\,HF(f).
\]
恒等也被保留。对象和态射都不改变的函子记为 \(\operatorname{Id}_{\mathcal C}\)，复合结合律及恒等律逐对象、逐态射成立。

\begin{proposition}\label{b4:prop:vertical-composition}
设 \(\mathcal C,\mathcal D\) 为范畴，\(F,G,H:\mathcal C\rightarrow\mathcal D\) 为函子，
\(\eta:F\Rightarrow G\)、\(\theta:G\Rightarrow H\) 为自然变换。则
\[
(\theta\eta)_X=\theta_X\eta_X
\]
定义自然变换 \(\theta\eta:F\Rightarrow H\)，称为 \(\eta\) 与 \(\theta\) 的\term[chuizhifuhe]{垂直复合}[vertical composition]，即逐分量复合。这种复合具有结合律；分量为 \(1_{F(X)}\) 的变换 \(1_F\) 为恒等变换。
\end{proposition}
\begin{proof}
对 \(f:X\rightarrow Y\)，依次使用两个自然性等式得到
\[
H(f)\theta_X\eta_X
=\theta_YG(f)\eta_X
=\theta_Y\eta_YF(f).
\]
所以复合自然。三个变换的复合在对象 \(X\) 处是三个 \(\mathcal D\) 态射的复合，结合律及恒等律直接来自 \(\mathcal D\)。
\end{proof}

\begin{definition}\label{b4:def:natural-isomorphism}
设 \(\mathcal C,\mathcal D\) 为范畴，\(F,G:\mathcal C\rightarrow\mathcal D\) 为函子，\(\eta:F\Rightarrow G\) 为自然变换。若存在自然变换 \(\theta:G\Rightarrow F\)，使
\(\theta\eta=1_F\)、\(\eta\theta=1_G\)，称 \(\eta\) 为\term[zirantonggou]{自然同构}[natural isomorphism]，记为 \(F\cong G\)。
\end{definition}

\begin{proposition}\label{b4:prop:natural-iso-components}
设 \(\mathcal C,\mathcal D\) 为范畴，\(F,G:\mathcal C\rightarrow\mathcal D\) 为函子，\(\eta:F\Rightarrow G\) 为自然变换。则 \(\eta\) 为自然同构，当且仅当每个分量 \(\eta_X\)（\(X\in\mathcal C\)）都是 \(\mathcal D\) 中的同构。
\end{proposition}
\begin{proof}
自然逆的各分量显然是分量逆。反过来，假设每个分量可逆。由
\(G(f)\eta_X=\eta_YF(f)\)，左右分别复合相应逆态射，得到
\[
F(f)\eta_X^{-1}=\eta_Y^{-1}G(f).
\]
因此这些逆态射组成一个自然变换 \(G\Rightarrow F\)；逐对象复合为恒等，故是自然逆。逆态射唯一，所以这里不需要为每个对象另作一次选择。
\end{proof}

\begin{example}\label{b4:exm:objectwise-iso-not-natural}
设 \(k\) 为域，\(J\) 为只有对象 \(0,1\) 及一条非恒等态射 \(a:0\rightarrow1\) 的范畴。定义函子 \(F,G:J\rightarrow\mathbf{Vect}_k\)，令
\[
F(0)=F(1)=G(0)=G(1)=k,\qquad
F(a)=1_k,\qquad G(a)=0,
\]
并令两函子把恒等态射送到 \(1_k\)。它们在每个对象处的取值都相同，但若 \(\eta:F\Rightarrow G\) 为自然变换，对 \(a\) 的自然性便给出
\[
0\eta_0=\eta_1 1_k,
\]
所以 \(\eta_1=0\)，不可能是同构。因此 \(F\) 与 \(G\) 不自然同构。函子在对象处的取值既不能决定其态射作用，逐对象存在同构也不能保证存在自然同构。
\end{example}

自然变换还可以与函子在左右两侧复合。若 \(\eta:F\Rightarrow G\)，
\(H:\mathcal D\rightarrow\mathcal E\)，则
\((H\eta)_X=H(\eta_X)\) 给出 \(HF\Rightarrow HG\)，因为
\[
HG(f)\,H(\eta_X)
=H\bigl(G(f)\eta_X\bigr)
=H\bigl(\eta_YF(f)\bigr).
\]
若 \(K:\mathcal B\rightarrow\mathcal C\)，则
\((\eta K)_B=\eta_{K(B)}\) 给出 \(FK\Rightarrow GK\)，自然性就是 \(\eta\) 对箭头 \(K(u)\) 的自然性。

\begin{proposition}\label{b4:prop:horizontal-composition}
设 \(\mathcal C,\mathcal D,\mathcal E\) 为范畴，\(F,G:\mathcal C\rightarrow\mathcal D\)、\(H,L:\mathcal D\rightarrow\mathcal E\) 为函子，
\(\eta:F\Rightarrow G\)、\(\alpha:H\Rightarrow L\) 为自然变换。从 \(HF\) 到 \(LG\)，可以先改变内侧函子，再改变外侧函子，也可以按相反次序复合：
\[
\begin{aligned}
HF&\xRightarrow{H\eta}HG\xRightarrow{\alpha G}LG,\\
HF&\xRightarrow{\alpha F}LF\xRightarrow{L\eta}LG.
\end{aligned}
\]
这两条路径给出同一个自然变换，其分量为
\[
(\alpha*\eta)_X
=\alpha_{G(X)}H(\eta_X)
=L(\eta_X)\alpha_{F(X)}
\]
称这个自然变换为 \(\eta\) 与 \(\alpha\) 的\term[shuipingfuhe]{水平复合}[horizontal composition]，也称横向复合。
\end{proposition}
\begin{proof}
两个公式相等，正是 \(\alpha\) 对 \(\mathcal D\) 中态射
\(\eta_X:F(X)\rightarrow G(X)\) 的自然性。第一个公式表示自然变换
\(H\eta:HF\Rightarrow HG\) 与 \(\alpha G:HG\Rightarrow LG\) 的逐分量复合，故由\cref{b4:prop:vertical-composition}自然。
\[
\begin{tikzcd}[column sep=large,row sep=large]
H(F(X))\arrow[r,"H(\eta_X)"]\arrow[d,"\alpha_{F(X)}"']&
H(G(X))\arrow[d,"\alpha_{G(X)}"]\\
L(F(X))\arrow[r,"L(\eta_X)"']&L(G(X)).
\end{tikzcd}
\]
\end{proof}

这项公式用于比较复合函子。把范畴看作对象、函子看作1态射、自然变换看作2态射，两种复合共同形成严格2范畴的例子。它们之间的交换律及完整结构见\secref{b4:sec:F01}。

\subsection{忘却结构与附加结构}
\label{b4:subsec:0102:03}

忘却函子保留底层映射，只忘掉一部分对象结构。例如
\[
R\text{-}\mathbf{Mod}\longrightarrow\mathbf{Ab}
\longrightarrow\mathbf{Set}
\]
先忘掉 \(R\) 的作用，再忘掉加法。它们确实是函子，因为恒等与复合的底函数都没有改变。反方向则不能把一个任意集合直接宣称为模；必须构造新的元素及运算。

\begin{example}\label{b4:ex:free-module-functor}
固定含幺环 \(R\)。第二卷第3.1节构造的自由左模
\[
R^{(X)}=\left\{\sum_{x\in X}r_xe_x:
r_x\in R,\ \text{只有有限多个}\;r_x\ne0\right\}
\]
给出函子 \(F_R:\mathbf{Set}\rightarrow R\text{-}\mathbf{Mod}\)。
对函数 \(u:X\rightarrow Y\)，规定
\[
F_R(u)\left(\sum_xr_xe_x\right)=\sum_xr_xe_{u(x)}.
\]
若不同 \(x\) 有同一个像，就把相应系数相加；支集有限保证和有意义。这个公式保持加法与左标量作用。又有
\(F_R(1_X)(e_x)=e_x\) 及
\(F_R(vu)(e_x)=e_{v(u(x))}=F_R(v)F_R(u)(e_x)\)，由基生成性可知它保持恒等与复合。
\end{example}
\symidx{\(R^{(X)}\)：以集合为基的自由左模}
\symidx{\(F_R\)：自由模函子}

设 \(U:R\text{-}\mathbf{Mod}\rightarrow\mathbf{Set}\) 为忘却函子。映射
\(\iota_X:X\rightarrow U F_R(X)\)、\(x\mapsto e_x\)，组成自然变换
\(\operatorname{Id}_{\mathbf{Set}}\Rightarrow UF_R\)，因为
\(UF_R(u)(e_x)=e_{u(x)}\)。另一个方向有线性映射
\[
\epsilon_M:F_RU(M)\longrightarrow M,\qquad
\sum_{m\in M}r_me_m\longmapsto\sum_mr_mm.
\]
对模同态 \(f:M\rightarrow N\)，在基元 \(e_m\) 上，
\(f\epsilon_M(e_m)=f(m)=\epsilon_NF_RU(f)(e_m)\)，所以这些映射也自然。这两族映射并非一对自然逆。虽然
\(U(\epsilon_M)\iota_{U(M)}=1_{U(M)}\)，但当 \(R\ne0\) 时，\(F_RU(R)\) 中对应零元素的基元 \(e_0\) 是非零向量，而 \(\epsilon_R(e_0)=0\)。因此 \(\epsilon_R\) 不单射，不能是同构；\(e_0\) 与自由模中的零向量必须区别。第3章会说明这两族映射怎样组成自由与忘却之间的伴随。

\ProseParagraphBreak
构造是否能够成为函子，还取决于允许哪些态射。例如群的中心 \(Z(G)\) 对同构有相容的映射，但任意同态
\(f:G\rightarrow H\) 未必把 \(Z(G)\) 送入 \(Z(H)\)。取
\(G=\{1,(12)\}\subset S_3=H\)，则 \(G\) 交换，中心为自身，而 \(S_3\) 的中心平凡；包含映射便不能限制为 \(Z(G)\rightarrow Z(H)\)。问题不是中心不存在，而是所提议的态射规则不成立。
